\section{The complete integer two-scale identity}
\label{clc:sec:two}

The general table construction has a particularly short closed form
for scales $1$ and $m$, where $m$ is any positive integer. Write
\begin{equation}\label{clc:eq:pair-data}
 \sigma_m=(-1)^{m+1},\qquad
 c_r^{(m)}=(-1)^{r+1}\csc\frac{\pi r}{m}\quad(1\le r<m).
\end{equation}
When $m=1$ all sums over $r$ below are empty.
For $\Rea A>-1$, define the \emph{unnormalized} integral
\begin{equation}\label{clc:eq:I-def}
 I_m(s,t;A)=\int_0^\infty
 F_{1+A,s}(y)F_{1+A/m,t}(y^{-m})\frac{\dd y}{y}.
\end{equation}
The normalized quantity is
$D_m(W;A)=\CC_{(1,m),A}(W;0)=m^{-t}I_m(s,t;A)$.

\begin{theorem}[All integer scale ratios and all complex orders]
\label{clc:thm:pair}
For $s,t\in\C$, $W=s+t$, and $\Rea A>-1$,
\begin{align}
 I_m(s,t;A)=m^{-s-1}\biggl[&
 \pi\sum_{r=1}^{m-1}c_r^{(m)}
       \Phi\left(\sigma_m,W,\frac{A+r}{m}\right)\notag\\
 &+\sigma_m W\Phi\left(\sigma_m,W+1,1+\frac A m\right)\biggr].
 \label{clc:eq:pair}
\end{align}
The integral is ordinary, absolutely convergent, and entire in $s,t$.
For odd $m$, the last term is $\EE_{1+A/m}(W)$, and the simple Hurwitz
sum is grouped into differences to remove its pole at $W=1$.
\end{theorem}
\begin{proof}
Use $Q=m$. The poles at $r=1,\ldots,m-1$ are simple, with
\[
 a_{r,1}=(-1)^r\frac\pi m\csc\frac{\pi r}{m}.
\]
At $\rho=m$ both factors are polar, and
$a_{m,2}=(-1)^{m+1}=\sigma_m$, $a_{m,1}=0$.
Substitution in \eqref{clc:eq:central-raw} gives the normalized formula
with prefactor $m^{-W-1}$. Multiplication by $m^t$ gives
\eqref{clc:eq:pair}. Entire dependence and absolute convergence were
proved before the residue calculation. For odd $m$,
$c_{m-r}^{(m)}=-c_r^{(m)}$, so $\sum_rc_r^{(m)}=0$ and the
simple-pole cancellation is explicit.
\end{proof}

\subsection{Scales two and three}
At $A=0$ the functions in the integrand are ordinary polylogarithms.
For $m=2$, the two reductions
$\eta(W,1/2)=2^W\beta(W)$ and $\eta(W,1)=\eta(W)$ give
\eqref{clc:eq:intro-beta}. For $m=3$, let $\chi_{-3}$ be the real
nonprincipal character modulo three and set
\[
 L(W,\chi_{-3})=3^{-W}
       [\zeta(W,1/3)-\zeta(W,2/3)].
\]
Then
\begin{equation}\label{clc:eq:pair3}
 \boxed{\displaystyle
 \int_0^\infty\Li_s(-y)\Li_t(-y^{-3})\frac{\dd y}{y}
 =\frac{2\pi\,3^{t-1}}{\sqrt3}L(W,\chi_{-3})
    +3^{-s-1}\EE_1(W).}
\end{equation}
For $s=t=1$ this yields
\begin{equation}\label{clc:eq:log3}
 \int_0^\infty\log(1+y)\log(1+y^{-3})\frac{\dd y}{y}
 =\frac{2\pi}{\sqrt3}L(2,\chi_{-3})+\frac29\zeta(3).
\end{equation}
The first two numerical values are
\begin{align*}
 I_2(1,1;0)&=2.426819442923313294752678355954754053578\ldots,\\
 I_3(1,1;0)&=3.101375395830302048787723234392638\ldots.
\end{align*}
They were checked against direct real-line logarithmic quadrature;
they are diagnostics, not the proof.

\subsection{Zero-order resonance and an explicit normalization obstruction}
At zero orders one obtains
\begin{equation}\label{clc:eq:zero-examples}
 I_2(0,0;0)=\frac\pi4,\qquad
 I_3(0,0;0)=\frac13+\frac{2\pi}{9\sqrt3}.
\end{equation}
The $1/3$ in the second expression is the value
$3^{-1}\EE_1(0)$. Substituting $W=0$ in a separated factor
$W\zeta(W+1)$ and declaring it zero would lose this term.

Even at the smallest nonzero total order the raw unequal-scale
integral fails pure order addition:
\begin{align}
 I_2(1,0;0)&=\frac{5\pi^2}{48},&
 I_2(0,1;0)&=\frac{5\pi^2}{24}.
 \label{clc:eq:swap}
\end{align}
For instance, use $\beta(1)=\pi/4$ and $\eta(2)=\pi^2/12$ in
\eqref{clc:eq:intro-beta}. These two values have the same total order
but differ by a factor of two. The normalized integrals agree, as
Theorem~\ref{clc:thm:addition} requires. This is an exact counterexample
to a tempting extension of the old equal-scale statement, not an error
attributed to that statement itself.

For $m=1$, Theorem~\ref{clc:thm:pair} reduces to the preceding report's
identity $I_1(s,t;A)=\EE_{1+A}(s+t)$. This limiting consistency is useful,
but the equal-scale case is prior work, not a new claim.
